\begin{proposition}[The dominant half]\label{prop:dominant-half}
Under the hypotheses of Proposition~\ref{prop:all-bias-modes}, suppose
$q\ge1/2$. Then $P_n(n-k)\le P_n(k)$ for $2k\ge n$, and the right half
has no strict valley:
\[
 P_n(j)\ge\min\{P_n(i),P_n(k)\}
 \qquad\text{if }i<j<k\text{ and }2i+1\ge n.
\]
Equivalently, its first differences contain no negative value followed
by a positive one. Reflection gives the corresponding left-half result
when $q\le1/2$.
\end{proposition}

\begin{proof}
Temporarily display the bias in $P_{n,q}$ and put
\[
 A_n(k)=P_{n,q}(k)-P_{n,1-q}(k)
       =P_{n,q}(k)-P_{n,q}(n-k).
\]
The forward equation gives
$A_{n+1}(k)=(1-r_n(k))A_n(k)+r_n(k-1)A_n(k-1)$.
Induction yields $A_n(k)\ge0$ for $2k\ge n$: both parents lie in the
previous upper half unless $k$ is the new centre, where $A_{n+1}(k)=0$
by reflection; terms outside the support vanish. In particular, the
central first difference of an odd row is nonnegative.

Restrict the augmented row $Z_t$ in the preceding proof to spatial
indices $2k+1\ge t$. At $t=1$ its pattern is contained in $(+,-)$:
the central entry is $2q-1\ge0$, and the subsequent entries are the
future right boundary values $r_s^*(q)=-\ell_s(1-q)$. Since
$1-q\le1/2$, their pattern is $(+,-)$ by
\eqref{eq:bias-edge-increments}. The positive two-parent recurrence
preserves absence of a negative entry followed by a positive one.
The sole possible parent outside the previous half occurs at the new
central difference of an odd row. That entry is nonnegative, so it
cannot begin the forbidden pair. More explicitly, if a negative entry
at index $i$ of row $t+1$ satisfies $2i+1\ge t+1$, then $2i>t$;
both its parents satisfy $2k+1\ge t$. Parents of a later positive
entry follow them strictly, as in the preceding proof. Induction
therefore excludes the pair in every restricted augmented row.
A strict valley would contain such two adjacent differences by
telescoping, which proves the claim.
\end{proof}
